\documentclass[10pt]{article}
\usepackage[T1]{fontenc}
\usepackage[utf8]{inputenc}
\usepackage[french]{babel}
\usepackage{amsmath,amssymb,mathrsfs,mathpazo}
\usepackage{geometry}
\geometry{a4paper,margin=2cm}
\usepackage{graphicx,array,tabularx,tabularray,colortbl,multirow,multicol,booktabs,enumitem,pifont,ulem,fancybox,diagbox,cancel}
\usepackage{pstricks,pstricks-add,pst-plot,pst-tree,pst-node,pst-text,pst-func,pst-math,pst-bspline,pst-3dplot,pst-eucl,pst-solides3d}
\let\pscalculate\undefined
\usepackage{pst-bezier}
\usepackage{tikz,pgfplots,tkz-tab}
\usetikzlibrary{arrows,arrows.meta,patterns,calc,positioning,trees,decorations.markings,shapes.geometric,shapes.misc,shapes.arrows,intersections,decorations.pathmorphing,decorations.pathreplacing,backgrounds,shadows,babel}
\pgfplotsset{compat=1.18}
\usepgfplotslibrary{fillbetween}
\usepackage[np]{numprint}
\usepackage{listings,makecell,pas-tableur,twoopt}
\let\barre\undefined
\usepackage[pstricks]{tablvar}
\definecolor{darkgreen}{RGB}{0,100,0}
\providecommand{\red}{\color{red}}
\providecommand{\blue}{\color{blue}}
\providecommand{\green}{\color{green}}
\providecommand{\black}{\color{black}}
\providecommand{\white}{\color{white}}
\providecommand{\gray}{\color{gray}}
\providecommand{\cyan}{\color{cyan}}
\providecommand{\magenta}{\color{magenta}}
\providecommand{\yellow}{\color{yellow}}
\providecommand{\violet}{\color{violet}}
\providecommand{\orange}{\color{orange}}
\providecommand{\euro}{\texteuro}
\providecommand{\EUR}{\texteuro}
\providecommand{\ieme}{\textsuperscript{e}}
\providecommand{\degres}{\ensuremath{^\circ}}
\providecommand{\degree}{\ensuremath{^\circ}}
\providecommand{\vv}[1]{\overrightarrow{#1}}
\providecommand{\Coord}[1]{#1}
\providecommand{\decosix}{}
\providecommand{\starredbullet}{\ensuremath{\bullet}}
\pagestyle{empty}
\setlength{\parindent}{0pt}
\usepackage{adjustbox}\DeclareUnicodeCharacter{2236}{:}\DeclareUnicodeCharacter{2190}{\ensuremath{\leftarrow}}
\providecommand{\R}{\mathbb{R}}
\providecommand{\N}{\mathbb{N}}
\providecommand{\Z}{\mathbb{Z}}
\providecommand{\Q}{\mathbb{Q}}
\providecommand{\D}{\mathbb{D}}
\providecommand{\C}{\mathbb{C}}
\providecommand{\e}{{\mathrm{e}}}
\providecommand{\vect}[1]{\overrightarrow{#1}}
\providecommand{\vectt}[1]{\overrightarrow{\mathrm{#1}}}
\providecommand{\vvt}[1]{\overrightarrow{\mathrm{#1}}}
\providecommand{\np}[2][]{{#2~#1}}
\providecommand{\vv}[1]{{\overrightarrow{#1}}}
\renewcommand{\d}{{\mathrm{d}}}
\providecommand{\degres}{{^\circ}}
\providecommand{\degre}{{^\circ}}
\providecommand{\blue}{\color{blue}}
\providecommand{\red}{\color{red}}
\providecommand{\enskip}{\quad}
\providecommand{\llbracket}{\lbrack\!\lbrack}
\providecommand{\rrbracket}{\rbrack\!\rbrack}
\providecommand{\euro}{{\text{€}}}
\providecommand{\no}{{\text{n°}}}
\providecommand{\interval}[2]{{[#1;#2]}}
\providecommand{\intervFO}[2]{{[#1;#2[}}
\providecommand{\dotfill}{\ldots}

\providecommand{\ds}{\displaystyle}
\providecommand{\textding}[1]{#1}
\providecommand{\Oij}{\(\mathrm{O};\vec{\imath},\vec{\jmath}\)}
\providecommand{\Oijk}{\(\mathrm{O};\vec{\imath},\vec{\jmath},\vec{k}\)}
\providecommand{\Ouv}{\(\mathrm{O};\vec{u},\vec{v}\)}
\providecommand{\intervFF}[2]{\ensuremath{[#1;#2]}}
\renewcommand{\euro}{\ifmmode\text{\texteuro}\else\texteuro\fi}
\DeclareUnicodeCharacter{00B0}{\ifmmode^{\circ}\else\textdegree\fi}
\DeclareUnicodeCharacter{20AC}{\ifmmode\text{\texteuro}\else\texteuro\fi}
% Énoncés : Éducation nationale. Transcription et corrigés : APMEP, Denis Vergès et contributeurs. Découpage et classement : CoopMaths, Sébastien Lozano, Éric Elter et contributeurs. Adaptation de la présentation : Maths Explorer.
% CC BY-SA : https://coopmaths.fr/www/about/
% https://coopmaths.fr/alea/static/bac/2021/tex/bac_2021_06_etrangers_10.tex
\begin{document}
\noindent\textit{APMEP / CoopMaths. CC BY-SA. Présentation adaptée par Maths Explorer.}\par\medskip
\medskip

\textbf{Partie A : Détermination d'une fonction $f$ et résolution d'une équation différentielle}

\medskip

On considère la fonction $f$ définie sur $\R$ par :

\[f(x) = \e^x+ ax + b\e^{-x}\]

où $a$ et $b$ sont des nombres réels que l'on propose de déterminer dans cette partie.

Dans le plan muni d'un repère d'origine O, on a représenté ci-dessous la courbe $\mathcal{C}$, représentant la fonction $f$, et la tangente $(T)$ à la courbe $\mathcal{C}$ au point d'abscisse $0$.

\begin{center}
\psset{unit=1.25cm}
\begin{pspicture*}(-2.2,-0.6)(3.2,6.6)
\psgrid[gridlabels=0pt,subgriddiv=5,gridwidth=0.25pt,subgridwidth=0.1pt]
\psaxes[linewidth=1.25pt,labelFontSize=\scriptstyle]{->}(0,0)(-2.2,-0.4)(3.2,6.6)
\psplot[plotpoints=2000,linewidth=1.25pt,linecolor=red]{-2}{2.5}{2.71828 x exp x sub 2 2.71828 x exp div add} \uput[r](2,5){\red $\mathcal{C}$}
\psplot[linewidth=1.25pt]{-2}{2}{3 2 x mul sub}
\end{pspicture*}
\end{center}

\smallskip

\begin{enumerate}
\item Par lecture graphique, donner les valeurs de $f(0)$ et de $f’(0)$.
\item En utilisant l'expression de la fonction $f$, exprimer $f(0)$ en fonction de $b$ et en déduire la valeur de $b$.
\item On admet que la fonction $f$ est dérivable sur $\R$ et on note $f'$ sa fonction dérivée. 
	\begin{enumerate}
		\item Donner, pour tout réel $x$, l'expression de $f'(x)$.
		\item Exprimer $f'(0)$ en fonction de $a$.
		\item En utilisant les questions précédentes, déterminer $a$, puis en déduire l'expression de $f(x)$.
	\end{enumerate}
\item On considère l'équation différentielle :

\[(E):\quad  y' +y =2\e^x - x - 1\]

	\begin{enumerate}
		\item Vérifier que la fonction $g$ définie sur $\R$ par :
		
\[g(x) = \e^x - x + 2\e^{-x}.\]

est solution de l'équation $(E)$.
		\item Résoudre l'équation différentielle $y' + y = 0$.
		\item En déduire toutes les solutions de l'équation $(E)$. 
est solution de l'équation (E).
	\end{enumerate}
\end{enumerate}

\bigskip

\textbf{Partie B : Étude de la fonction $g$ sur $[1~;~+\infty[$}

\medskip

\begin{enumerate}
\item Vérifier que pour tout réel $x$, on a :

\[\e^{2x} - \e^x - 2 = \left(\e^x - 2\right)\left(\e^x + 1\right)\]

\item En déduire une expression factorisée de $g'(x)$, pour tout réel $x$.
\item On admettra que, pour tout $x \in  [1~;~ +\infty[$,\, $\e^x - 2 > 0$.

 Étudier le sens de variation de la fonction $g$ sur $[1~;~ +\infty[$.
\end{enumerate}

\end{document}
